\begin{abstract}
Agent tasks cross model, tool and Web services, and their encrypted traffic exposes information at several execution stages. We systematize fifteen Agent studies (forty-four settings) using twenty-two traffic-analysis anchors. The synthesis connects two evaluation requirements: follow a task across every service view available to its attacker, and measure what executed cover does to the original business task. None of the four mapped query/workflow privacy settings documents both an adapted attacker and matched task utility. We then examine cover on the same MCP executions from two fixed-weight LLM planners. A size-and-destination attack recognizes 61--63 of 120 covered tasks for Qwen and 72--82 for Llama. Adding client-receipt timing enables a two-stage attack: detect controller-inserted decoys, remove them, then classify the remaining calls. Recognition rises to 111--114 and 107--115, respectively; all ten planned secondary contrasts survive correction. Simply concatenating timing features gives no improvement. The controller dispatches decoys without an intervening model decision, leaving a timing distinction despite padded messages. These results link the observer model to a concrete cover-removal mechanism, alongside the separate state costs of adverse writes. We contribute access-aware attack and defense maps, a classical-to-Agent transfer matrix and an evaluation procedure that binds residual inference to task, state and resource outcomes.
\end{abstract}
